% !TEX encoding = UTF-8 Unicode
% !TeX program = xelatex
% !TEX TS-program = XeLaTeX

\documentclass[12pt, a4paper]{article}
%\documentclass[12pt, a5paper]{article}
\usepackage{mathtools}
\usepackage{unicode-math}
\usepackage{comment}
\usepackage{wrapfig}
\usepackage{tikz}
\usetikzlibrary{arrows, decorations.markings}
\usepackage{subcaption}
\usepackage{subfig}


\usepackage[pass]{geometry}
%\geometry{legalpaper, landscape, margin=1in}
\usepackage{fullpage}
\usepackage{setspace}
\usepackage[localise]{xepersian}
\setstretch{1.255}
\settextfont[Scale=1]{XB Niloofar}
\setlatintextfont[Scale=1]{XB Niloofar}%{Junicode} %Junicode
\setdigitfont{XB Niloofar}%{XB Zar}

\setlength{\parindent}{0ex}

\definecolor{answercolor}{HTML}{0080ff}
%\newcommand{\پاسخ}[2]{#1}
\newcommand{\پاسخ}[2]{\color{answercolor}\textbf{پاسخ: }#2\color{black}\vspace{-0.3cm}}

\begin{document}
	\tikzset{
		treenode/.style = {align=center, inner sep=0pt, text centered,
			font=\sffamily},
		arn_n/.style = {treenode, circle, white, font=\sffamily\bfseries, draw=black,
			fill=black, text width=1.5em},% arbre rouge noir, noeud noir
		arn_r/.style = {treenode, circle, red, draw=red, 
			text width=1.5em, very thick},% arbre rouge noir, noeud rouge
		arn_s/.style = {treenode, circle, black, draw=black, 
			text width=1.5em, very thick},% arbre rouge noir, noeud rouge
		arn_x/.style = {treenode, rectangle, draw=black,
			minimum width=0.5em, minimum height=0.5em, very thick}% arbre rouge noir, nil
	}
	\begin{tikzpicture}[remember picture,overlay]	\node[anchor=north,yshift=-1.5cm] at (current page.north) {\rl{به نام خدا}};
	\end{tikzpicture}%

\vspace{-1cm}
{\centering
	\textbf{کوئیز دوم درس ساختمان داده‌ها و الگوریتم‌ها}\\
مدرس: مهدی صفرنژاد       - ۲۰ اسفند ۹۷\\
\raggedleft\small
نام و نام خانوادگی: \\
}
\vspace{-0.5cm}
\پرش‌متوسط\پرخط‌ا\\
در برخی سؤالات اَعمالی را باید روی داده ساختار انجام دهید و نتیجه نهایی را نشان دهید. در برخی سؤالات فقط باید یک داده ساختار یا الگوریتم را نام ببرید. در سؤالات دیگر پاسخ کوتاه (حداکثر یک خط) مدنظر است.

\begin{wrapfigure}[2]{L}{0.4\textwidth}
\vspace{-1cm}
	\begin{center}
		\begin{tikzpicture}[-,>=stealth',level/.style={sibling distance = 2cm/#1,
			level distance = 1cm}] 
		\node [arn_s] {20} 
		child{ node [arn_s] {10} 
			child{ node [arn_s] {6} 
			}
			child{ node [arn_s] {7}
			}
		}
		child { node [arn_s] {5}
%			child [grow=-105] { node [arn_s] {3} 
%			}
		}
		; 
		\end{tikzpicture}
	\end{center}
	\caption{هرم بیشینۀ سؤال ۱}
	\label{شکل:۱}
\end{wrapfigure}

\پاراگراف{سؤال ۱. }
عنصر ۲۵ را در هرم بیشینۀ شکل~\رجوع{شکل:۱} درج کنید و هرم پس از درج را رسم کنید.

\پاسخ{\vspace{3.5cm}}{
	
درج در شکل~\رجوع{شکل:۱جواب} نشان داده شده است.

\vspace{2.2cm}
\begin{figure}[h!]
	\begin{center}
		\begin{tikzpicture}[-,>=stealth',level/.style={sibling distance = 2cm/#1,
			level distance = 1cm}] 
		\node [arn_s] {20} 
child{ node [arn_s] {10} 
	child{ node [arn_s] {6} 
	}
	child{ node [arn_s] {7}
	}
}
child { node [arn_s] {5}
}
		; 
		
		\draw [decoration={markings,mark=at position 1 with
			{\arrow[scale=3,>=stealth]{>}}},postaction={decorate}] (2,-1) -- (4,-1);
		
		\node [arn_s] at (6,0) {25} 
child{ node [arn_s] {10} 
	child{ node [arn_s] {6} 
	}
	child{ node [arn_s] {7}
	}
}
child { node [arn_s] {20}
	child [grow=-105] { node [arn_s] {5} }
}
; 
		\end{tikzpicture}
	\end{center}
	\caption{پاسخ سؤال ۱}
	\label{شکل:۱جواب}
\end{figure}

}

\پاراگراف{سؤال ۲.}
آیا در هرم بیشینه می‌توان اَعمال درج، پیدا کردن عنصر بیشینه و پیدا کردن عنصر کمینه را در $O(\log n)$ انجام داد؟

\پاسخ{\vspace{0.5cm}}{
خیر، پیدا کردن عنصر کمینه در هرم بیشینه در زمان $O(\log n)$ میسر نیست.
	
}

\پاراگراف{سؤال ۳.}
اگر یک دوران چپ‌گرد روی یک رأس و سپس یک دوران راست‌گرد روی \زیرخط{همان} رأس یک درخت دودویی جست‌وجو انجام دهیم، ساختار درخت به حالت اول بازمی‌گردد؟

\پاسخ{\vspace{0.7cm}}{
خیر. مثال نقض در شکل~\رجوع{شکل:۳جواب} نشان داده شده است. در پاسخ اشاره به این که باید دوران دوم روی رأس دیگری انجام شود تا به حالت اول برگردد نیز کفایت می‌کند.

\vspace{0.5cm}
\begin{figure}[h!]
	\begin{center}
		\begin{tikzpicture}[-,>=stealth',level/.style={sibling distance = 2cm/#1,
			level distance = 1cm}] 
		\node [arn_s] at (0,0) {x} 
		child{ node [arn_s] {a} 
		}
		child{ node [arn_s] {y}
			child{ node [arn_s] {b} 
			}
			child{ node [arn_s] {c}
			}
		}
		; 

		\draw [decoration={markings,mark=at position 1 with
			{\arrow[scale=3,>=stealth]{>}}},postaction={decorate}] (2,-1) -- (4,-1);
		
		\node [arn_s] at (6,0) {y}
		child{ node [arn_s] {x} 
			child{ node [arn_s] {a} 
			}
			child{ node [arn_s] {b}
			}
		}
		child{ node [arn_s] {c}
		}
		; 
		
		\draw [decoration={markings,mark=at position 1 with
			{\arrow[scale=3,>=stealth]{>}}},postaction={decorate}] (8,-1) -- (10,-1);

		\node [arn_s] at (12,0) {y}
		child{ node [arn_s] {a} 
			child [grow=south east] { node [arn_s] {x}
				child{ node [arn_s] {b}
				}
			}
		}
		child{ node [arn_s] {c}
		}
		; 
		\end{tikzpicture}
	\end{center}
	\caption{پاسخ سؤال ۳}
	\label{شکل:۳جواب}
\end{figure}

}
	
\پاراگراف{سؤال ۴.}
آیا اگر ساختار یک درخت دودویی جست‌وجو و مجموعه عناصر ذخیره شده در آن را بدانیم، می‌توانیم عنصر ذخیره شده در هر رأس را مشخص کنیم؟

\پاسخ{\vspace{0.5cm}}{
بله، با توجه به این که در پیمایش میان‌ترتیب عناصر به صورت صعودی هستند، می‌توان عنصر مربوط به هر رأس را پیدا کرد.
	
}

\پاراگراف{سؤال ۵.}
آیا در حین درج یک عنصر جدید در یک درخت دودویی جست‌وجو، از عنصر قبلی و بعدی آن (در صورت وجود) می‌گذریم؟

\پاسخ{\vspace{0.7cm}}{
بله، عنصر بعدی یک رأس یا کوچکترین عنصر زیردرخت سمت راست است و یا در اجداد آن است. با توجه به این که رأسی که تازه درج شده زیردرختی ندارد، عنصر بعدی آن در اجداد آن است و در هنگام درج از آن می‌گذریم. همین استدلال برای عنصر قبلی برقرار است.
	
}

\begin{wrapfigure}[4]{l}{0.5\textwidth}
\vspace{-1.5cm}
	\begin{center}
		\begin{tikzpicture}[-,>=stealth',level/.style={sibling distance = 5cm/#1,
			level distance = 1cm}] 
		\node [arn_s] {10}
		child{ node [arn_s] {5} 
			child{ node [arn_x] {}}
			child{ node [arn_s] {6}
				child{ node [arn_x] {}}
				child{ node [arn_x] {}}
			}
		}
		child{ node [arn_s] {20}
			child{ node [arn_s] {15} 
				child{ node [arn_x] {}}
				child{ node [arn_s] {17}
					child{ node [arn_x] {}}
					child{ node [arn_x] {}}
				}
			}
			child{ node [arn_s] {25}
				child{ node [arn_x] {}}
				child{ node [arn_x] {}}
			}
		}
		; 
		\end{tikzpicture}
	\end{center}
	\caption{درخت مربوط به سؤال ۶}
	\label{شکل:۶}
\end{wrapfigure}

\پاراگراف{سؤال ۶.}
عنصر ۱۶ را در درخت ‌اِی‌وی‌اِل شکل~\رجوع{شکل:۶} درج کنید.

\پاسخ{\vspace{2.5cm}}{
درخت نهایی در شکل~\رجوع{شکل:۶جواب}	نشان داده شده است. یک دوران راست روی ۱۷ و یک دوران چپ‌گرد روی ۱۵ انجام می‌شود.
	
	\begin{figure}[h!]%{c}{0.5\textwidth}
		\vspace{2.5cm}
		\begin{center}
			\begin{tikzpicture}[-,>=stealth',level/.style={sibling distance = 5cm/#1,
				level distance = 1cm}] 
			\node [arn_s] {10}
			child{ node [arn_s] {5} 
				child{ node [arn_x] {}}
				child{ node [arn_s] {6}
					child{ node [arn_x] {}}
					child{ node [arn_x] {}}
				}
			}
			child{ node [arn_s] {20}
				child{ node [arn_s] {16} 
					child{ node [arn_s] {15}
						child{ node [arn_x] {}}
						child{ node [arn_x] {}}					
					}
					child{ node [arn_s] {17}
						child{ node [arn_x] {}}
						child{ node [arn_x] {}}
					}
				}
				child{ node [arn_s] {25}
					child{ node [arn_x] {}}
					child{ node [arn_x] {}}
				}
			}
			; 
			\end{tikzpicture}
		\end{center}
		\caption{پاسخ سؤال ۶}
		\label{شکل:۶جواب}
	\end{figure}

}

\begin{figure}[h!]%[5]{L}{0.4\textwidth}
	%	\vspace{-2cm}
	\begin{center}
		\begin{tikzpicture}[-,>=stealth',level/.style={sibling distance = 5cm/#1,
			level distance = 1cm}] 
		\node [arn_s] {5}
		child{ node [arn_s] {1} 
			child{ node [arn_x] {}}
			child{ node [arn_x] {}}
		}
		child{ node [arn_s] {15}
			child{ node [arn_s] {10} 
				child{ node [arn_s] {7}
					child{ node [arn_x] {}}
					child{ node [arn_x] {}}
				}
				child{ node [arn_s] {13}
					child{ node [arn_x] {}}
					child{ node [arn_x] {}}
				}
			}
			child{ node [arn_s] {20}
				child{ node [arn_x] {}}
				child{ node [arn_x] {}}
			}
		}
		; 
		\end{tikzpicture}
	\end{center}
	\caption{درخت مربوط به سؤال ۷ و ۸}
	\label{شکل:۷و۸}
\end{figure}

\پاراگراف{سؤال ۷.}
رأس‌های درخت شکل~\رجوع{شکل:۷و۸} را با رنگ‌های قرمز و سیاه طوری رنگ کنید که خاصیت‌های درخت قرمز-سیاه را داشته باشد.
	
\پاسخ{\vspace{4.5cm}}{
درخت رنگ‌آمیزی شده در شکل~\رجوع{شکل:۷} نشان داده شده است.

}

\پاراگراف{سؤال ۸.}
عنصر ۱۲ را در درخت قرمز-سیاهی که رنگ کردید درج کنید.

\پاسخ{\vspace{5cm}}{
	در شکل~\رجوع{شکل:۸} درخت نهایی نشان داده شده است. در واقع پس از عمل درج، یک بار تغییر رنگ، یک دوران راست‌گرد روی رأس ۱۵ و سپس یک دوران چپ‌گرد روی رأس ۵ داریم. اگر در پاسخ، رنگ رئوس ۵ و ۱۵ سیاه درنظر گرفته شود نیز با اغماض قابل قبول است.

\vspace{0.5cm}
%\begin{wrapfigure}{l}{0.5\textwidth}
\begin{figure}[h!]
%\vspace{2cm}
%	\begin{center}
\centering
\begin{subfigure}{0.45\textwidth}
		\begin{tikzpicture}[-,>=stealth',level/.style={sibling distance = 5cm/#1,
	level distance = 1cm}, scale=0.75] 
\node [arn_n] {5}
child{ node [arn_n] {1} 
	child{ node [arn_x] {}}
	child{ node [arn_x] {}}
}
child{ node [arn_r] {15}
	child{ node [arn_n] {10} 
		child{ node [arn_r] {7}
			child{ node [arn_x] {}}
			child{ node [arn_x] {}}
		}
		child{ node [arn_r] {13}
			child{ node [arn_x] {}}
			child{ node [arn_x] {}}
		}
	}
	child{ node [arn_n] {20}
		child{ node [arn_x] {}}
		child{ node [arn_x] {}}
	}
}
; 
\end{tikzpicture}
%	\end{center}
\caption{پاسخ سؤال ۷}
\label{شکل:۷}	
%\end{wrapfigure}
\end{subfigure}
\qquad
\begin{subfigure}{0.45\textwidth}
%\begin{wrapfigure}[10]{l}{0.5\textwidth}
%\begin{figure}[h!]
%	\begin{center}
		\begin{tikzpicture}[-,>=stealth',level/.style={sibling distance = 5cm/#1,
	level distance = 1cm}, scale=0.75] 
\node [arn_n] {10}
child{ node [arn_r] {5} 
	child{ node [arn_n] {1}
		child{ node [arn_x] {}}
		child{ node [arn_x] {}}
	}
	child{ node [arn_n] {7}
		child{ node [arn_x] {}}
		child{ node [arn_x] {}}
	}
}
child{ node [arn_r] {15}
	child{ node [arn_n] {13} 
		child{ node [arn_r] {12}
			child{ node [arn_x] {}}
			child{ node [arn_x] {}}
		}
		child{ node [arn_x] {}}			
	}
	child{ node [arn_n] {20}
		child{ node [arn_x] {}}
		child{ node [arn_x] {}}
	}
}
; 
\end{tikzpicture}
%	\end{center}
\caption{پاسخ سؤال ۸}
\label{شکل:۸}
	%\end{wrapfigure}
%\end{figure}
\end{subfigure}
\caption{پاسخ سؤال‌های ۷ و ۸}
\end{figure}

	
}

\پاراگراف{سؤال ۹.}
یک درخت جست‌وجوی متوازن که لزوماً دودویی نیست را نام ببرید.

\پاسخ{\vspace{0.5cm}}{
درخت ۲-۳، درخت ۲-۳-۴ و یا درخت بی

}

\پاراگراف{سؤال ۱۰.}
یک الگوریتم مرتب‌سازی نام ببرید که حافظۀ اضافۀ آن $O(1)$ و زمان اجرای آن در بدترین حالت $\Theta(n\log n)$ باشد.

\پاسخ{\vspace{0.5cm}}{
مرتب‌سازی هرمی

(مرتب‌سازی سریع با انتخاب هوشمندانه لولا در جلسات مربوط به امتحان نیست و حافظه اضافه آن $O(\log n)$ است، ولی با اغماض می‌توان آن را قبول کرد)
	
}

\begin{flushleft}
موفق باشید
\end{flushleft}

\end{document}